\documentclass[12pt,a4paper]{article}


% Encoding and language
\usepackage[utf8]{inputenc}
\usepackage[T1]{fontenc}
\usepackage[english]{babel}
\usepackage{enumitem}

% Mathematics
\usepackage{amsmath,amssymb,amsfonts}

% Graphics and figures
\usepackage{graphicx}
\graphicspath{{figuras_final/}}

% Tables
\usepackage{booktabs}
\usepackage{array}
\usepackage{longtable}
\usepackage{tabularx}
\usepackage{float}  % For exact positioning with [H]

% Links and references
\usepackage{hyperref}
\hypersetup{
    colorlinks=true,
    linkcolor=blue,
    citecolor=blue,
    urlcolor=blue,
    pdftitle={Reversal and Persistence of Gender Biases in GPT Models},
    pdfauthor={Otávio Oliveira Bopp}
}

% Bibliography
\usepackage{natbib}
\bibliographystyle{apalike}

% Margins
\usepackage[margin=2.5cm]{geometry}

% Spacing
\usepackage{setspace}
\onehalfspacing

\setcounter{tocdepth}{3}      % Show up to subsubsection in the table of contents
\setcounter{secnumdepth}{3}   % Number up to subsubsection
\renewcommand{\contentsname}{Table of Contents}

% For code
\usepackage{listings}
\lstset{
    basicstyle=\ttfamily\small,
    breaklines=true,
    frame=single
}

% For colored prompt boxes
\usepackage[most]{tcolorbox}
\usepackage{xcolor}

% Prompt box definition
\newtcolorbox{promptbox}[1][]{
  colback=gray!8,
  colframe=gray!60,
  fonttitle=\bfseries\small,
  title=#1,
  sharp corners,
  boxrule=0.8pt,
  left=8pt, right=8pt, top=6pt, bottom=6pt,
  fontupper=\small
}

% Box for prompts without title
\newtcolorbox{prompt}{
  colback=gray!8,
  colframe=gray!60,
  sharp corners,
  boxrule=0.8pt,
  left=8pt, right=8pt, top=6pt, bottom=6pt,
  fontupper=\small
}

% Metadata
\title{Reversal and Persistence of Gender Biases in GPT Models}
\author{Otávio Oliveira Bopp, Valdemar Pinho Neto}
\date{2025}

\begin{document}

\maketitle

\section{Appendix}

\subsection{Defining Bias}

% Hide subsubsections of 6.1 from the table of contents
\addtocontents{toc}{\protect\setcounter{tocdepth}{2}}



To understand the problem of bias in language models, we will
adopt an econometric perspective. Consider a hypothetical linear
regression, which predicts the probability of a candidate being a good
employee --- based on $i$ characteristics inferred from the résumé
and other information provided by the candidate, directly or indirectly
--- such that the candidates with the highest \textit{Y} are called to continue
the selection process (or even be hired without human
interference):

\begin{equation}
Y = \beta_0 + \sum_i (D_i \times \beta_i) + \varepsilon
\end{equation}

Where \textit{Y} is the dependent variable (the estimated ``quality'' of the
candidate), $D_i$ are dummy variables indicating the presence of each
characteristic $i$, and $\beta_i$ are the coefficients that measure the effect of
each characteristic on the dependent variable, $\varepsilon$ the error and $\beta_0$ the
intercept.

A regressor model is said to be \textit{unbiased} when the values of the estimated betas
($\hat{\beta}_i$) are equal to the real betas ($\tilde{\beta}_i$): that is, the
estimated effect is equal to the real effect of that explanatory variable on the
dependent variable. Formally:

\begin{equation}
E[\hat{\beta}_i] = \tilde{\beta}_i \quad \Rightarrow \quad \text{unbiased estimator}
\end{equation}

Among the $i$ characteristics, our hypothetical model will consider the
candidate's membership in the so-called ``protected classes''. Protected classes
are defined legally and may vary across different
jurisdictions, but aim to prevent discrimination against candidates based on
characteristics that have historically caused discrimination protected
class \citep{heymann2020}. Protected classes typically
include --- but are not limited to --- the candidate's race, gender
and sexual orientation. We will use the index $p$ to denote
protected characteristics, distinguishing them from the other characteristics
$i$.

A model will be called \textit{unbiased with respect to a protected characteristic}
when its estimated $\hat{\beta}_p$ is equal to zero --- which is taken
as a hypothesis as the real value $\tilde{\beta}_p$, with exceptions that will be discussed
in section 2.3. When we mention that a model was ``debiasing'',
we are saying that measures were taken so that the $\hat{\beta}_p$ of
protected characteristics approaches $\tilde{\beta}_p$.

The algorithms used in LLMs --- and Machine Learning models in
general --- do not follow a linear distribution, with well-defined
explanatory characteristics; rather, they find patterns in large
amounts of data, forming predictions without explicitly identifying
the explanatory variables. Even without defined explanatory variables or
a linear regression, we can infer that artificial intelligence models
take protected characteristics as part of their thousands
of latent variables --- given that they can infer these
characteristics from seemingly neutral information \citep{bai2024}. Thus, the previous definition can be expanded to include
these effects.

We break down the presence of a protected characteristic $p$ as a
summation of indicative variables of the ML model that allow inferring the
presence of that protected characteristic in the candidate. We denote this
aggregate effect as $\theta_p$:

\begin{equation}
\theta_p = \sum(\text{indicative variables of } p)
\end{equation}

The model will be called unbiased when the estimated value $\hat{\theta}_p$ is
equal to the real value $\tilde{\theta}_p$; which, by hypothesis, should be zero for the
majority of protected characteristics. When this relationship holds for
every $p$, the model is called unbiased.


\subsubsection{Proxy vs. Ideal --- Defining Discrimination}


The hypothesis that $\tilde{\beta}_p = 0$ assumes that, in an ideal world,
protected characteristics should not affect the quality of a
candidate. However, observable reality often diverges
from this ideal. The quotation from \citet{obermeyer2019} in the introduction illustrates
this problem: the health algorithm did not measure the need for medical
care (the characteristic that the model ideally tried to measure), but rather historical health
costs (an observable metric). Because unequal access to the health care system means that less was allocated to
Black patients, the algorithm systematically underestimated their
real needs.

We can call regressions with observable metrics --- such as health
costs, salary, or hiring rate --- \textit{proxy regressions}.
They indicate a characteristic that is not directly measurable. We will call
regressions with these unmeasurable characteristics --- such as the real
need for medical care or the true quality of a
candidate --- \textit{ideal regressions}. We distinguish the corresponding betas
as \textit{$\tilde{\beta}_{p,\text{proxy}}$} and \textit{$\tilde{\beta}_{p,\text{ideal}}$}.

From this distinction, we formally define \textit{discrimination} --- or
\textit{biased reality} --- as the difference between the real beta of a
proxy regression and the real beta of an ideal regression:

\begin{equation}
\text{Discrimination} = \tilde{\beta}_{p,\text{proxy}} - \tilde{\beta}_{p,\text{ideal}}
\end{equation}

By hypothesis, we maintain that, with rare exceptions, protected
characteristics should have \textit{$\tilde{\beta}_{p,\text{ideal}}$ = 0}. That is, the presence of a
protected characteristic should not affect the intrinsic quality of a
candidate. Therefore, outside these exceptions, discrimination emerges
when:

%\textit{$\tilde{\beta}_{p,\text{proxy}}$ $\neq$ 0 $\neq$ $\tilde{\beta}_{p,\text{ideal}}$}.
\[
\tilde{\beta}_{p,\text{proxy}} \neq 0 \neq \tilde{\beta}_{p,\text{ideal}}.
\]


\subsubsection{Exceptions to the Hypothesis $\tilde{\beta}_{p,\text{ideal}}$ = 0}


As mentioned earlier, it is not reasonable to believe that the real impact
of a protected characteristic on a candidate's quality is
\textit{always} zero. There are clear cases where $\tilde{\beta}_{p,\text{ideal}} \neq 0$. \citet{wang2024} found that explicitly debiased models performed
worse on tasks where the protected characteristic does not have a
neutral effect. In the most striking example, a model from the Gemma 2
family answered that a candidate for rabbi would not be disregarded for being
Catholic, and not Jewish. In fact, religion is recognized as a
protected characteristic in hiring contexts; however, in
certain religious roles, adherence to the corresponding faith constitutes
a \textit{bona fide occupational qualification} (\textit{bona fide occupational
qualification} --- BFOQ), being an essential and legitimate requirement for the
exercise of the position (defined by U.S. legislation in EEOC, CM-625).
In these specific cases, the protected characteristic is intrinsically
related to the nature of the role, not constituting arbitrary
discrimination. That is, $\tilde{\beta}_{p,\text{ideal}} \neq 0$.

Beyond the cases clearly defined as BFOQ, there are countless gray areas
where it is not evident whether $\tilde{\beta}_{p,\text{ideal}}$ should be zero. In another
example from \citet{wang2024}, the model from the Claude 3.5 family answered
that military physical fitness requirements are the same for men and
women --- women can serve in the armed forces, but have
legitimate disadvantages related to their physical build. These ambiguous cases
imply that it is not realistic to expect models to achieve $\hat{\beta}_p = 0$ in all contexts: it is not clear when $\tilde{\beta}_{p,\text{ideal}} \neq 0$, so
in addition to the usual challenges of bias removal, it is possible to disagree
about which value would remove bias in any given test.



% Restore depth to show 6.2.x in the table of contents
\addtocontents{toc}{\protect\setcounter{tocdepth}{3}}

\subsection{Full Regressions}

\subsubsection{Test 1: Characteristics in the Workplace Environment}

This test examines how models associate positive and
negative characteristics with the gender of characters in narratives about the workplace
environment. The dependent variable is is\_male (1 = male, 0 =
female/other). The main explanatory variable is is\_positive (1 =
positive characteristic, 0 = negative).

\textit{Note:} *** $p < 0,001$; ** $p < 0,01$; * $p < 0,05$. Base model for
dummies: gpt-3.5-turbo. Base order for shot order: MF (male
first).

\vspace{0.5cm}
\noindent\textbf{Regression 1: Chat Models --- Simple Specification}

\noindent N = 13.653 $|$ $R^2$ = 0,3803 $|$ k = 2

\begin{table}[H]
\centering
\begin{tabular}{lrrrrr}
\toprule
\textbf{Variable} & \textbf{Coef.} & \textbf{SE} & \textbf{t} & \textbf{p-value} & \textbf{Sig.} \\
\midrule
Intercept & 0,8490 & 0,0048 & 177,92 & $<$0,001 & *** \\
is\_positive & $-$0,6146 & 0,0067 & $-$91,73 & $<$0,001 & *** \\
\bottomrule
\end{tabular}
\end{table}

\noindent\textit{Interpretation: The negative coefficient of is\_positive ($-$0,61) indicates that
positive characteristics reduce the probability that the character is
male by 61 percentage points.}

\vspace{0.5cm}
\noindent\textbf{Regression 2: Chat Models --- With Language Control}

\noindent N = 13.653 $|$ $R^2$ = 0,4358 $|$ k = 3

\begin{table}[H]
\centering
\begin{tabular}{lrrrrr}
\toprule
\textbf{Variable} & \textbf{Coef.} & \textbf{SE} & \textbf{t} & \textbf{p-value} & \textbf{Sig.} \\
\midrule
Intercept & 0,7281 & 0,0056 & 130,02 & $<$0,001 & *** \\
is\_positive & $-$0,6142 & 0,0064 & $-$95,97 & $<$0,001 & *** \\
is\_portuguese & 0,2349 & 0,0064 & 36,70 & $<$0,001 & *** \\
\bottomrule
\end{tabular}
\end{table}

\noindent\textit{Interpretation: Controlling for language reveals that texts in Portuguese have
23pp more male characters than in English, but the effect of
is\_positive remains stable.}

\vspace{0.5cm}
\noindent\textbf{Regression 3: Chat Models --- Portuguese Only}

\noindent N = 7.019 $|$ $R^2$ = 0,3197 $|$ k = 2

\begin{table}[H]
\centering
\begin{tabular}{lrrrrr}
\toprule
\textbf{Variable} & \textbf{Coef.} & \textbf{SE} & \textbf{t} & \textbf{p-value} & \textbf{Sig.} \\
\midrule
Intercept & 0,9245 & 0,0066 & 140,08 & $<$0,001 & *** \\
is\_positive & $-$0,5372 & 0,0094 & $-$57,15 & $<$0,001 & *** \\
\bottomrule
\end{tabular}
\end{table}

\noindent\textit{Interpretation: In Portuguese, the pro-female bias is attenuated ($\beta$ = $-$0,54
vs.\ $-$0,61 overall), possibly due to grammatical gender marking.}

\vspace{0.5cm}
\noindent\textbf{Regression 4: Chat Models --- English Only}

\noindent N = 6.634 $|$ $R^2$ = 0,4966 $|$ k = 2

\begin{table}[H]
\centering
\begin{tabular}{lrrrrr}
\toprule
\textbf{Variable} & \textbf{Coef.} & \textbf{SE} & \textbf{t} & \textbf{p-value} & \textbf{Sig.} \\
\midrule
Intercept & 0,7689 & 0,0061 & 126,05 & $<$0,001 & *** \\
is\_positive & $-$0,6956 & 0,0086 & $-$80,88 & $<$0,001 & *** \\
\bottomrule
\end{tabular}
\end{table}
\noindent\textit{Interpretation: In English, the pro-female bias is stronger ($\beta$ = $-$0,70),
with a higher $R^2$ (0,50 vs.\ 0,32).}

\vspace{0.5cm}
\noindent\textbf{Regression 5: Chat Models --- Full Specification}

\noindent N = 13.653 $|$ $R^2$ = 0,5418 $|$ k = 15

\begin{table}[H]
\centering
\small
\begin{tabular}{lrrrrr}
\toprule
\textbf{Variable} & \textbf{Coef.} & \textbf{SE} & \textbf{t} & \textbf{p-value} & \textbf{Sig.} \\
\midrule
Intercept & 0,3234 & 0,0111 & 29,24 & $<$0,001 & *** \\
is\_positive & $-$0,6137 & 0,0058 & $-$106,23 & $<$0,001 & *** \\
is\_portuguese & 0,2428 & 0,0058 & 41,71 & $<$0,001 & *** \\
model\_gpt-4.1-2025-04-14 & 0,3583 & 0,0145 & 24,67 & $<$0,001 & *** \\
model\_gpt-4.1-mini-2025-04-14 & 0,4296 & 0,0145 & 29,58 & $<$0,001 & *** \\
model\_gpt-4.1-nano-2025-04-14 & 0,5875 & 0,0145 & 40,45 & $<$0,001 & *** \\
model\_gpt-4o-2024-08-06 & 0,4769 & 0,0145 & 32,83 & $<$0,001 & *** \\
model\_gpt-4o-mini & 0,5120 & 0,0145 & 35,26 & $<$0,001 & *** \\
model\_gpt-5-mini & 0,3861 & 0,0145 & 26,59 & $<$0,001 & *** \\
model\_gpt-5-nano & 0,2250 & 0,0145 & 15,49 & $<$0,001 & *** \\
model\_gpt-5.1-2025-11-13 & 0,4083 & 0,0145 & 28,12 & $<$0,001 & *** \\
model\_gpt-5.2-2025-12-11 & 0,2610 & 0,0165 & 15,82 & $<$0,001 & *** \\
model\_o3-2025-04-16 & 0,4204 & 0,0145 & 28,95 & $<$0,001 & *** \\
model\_o3-mini-2025-01-31 & 0,6694 & 0,0145 & 46,10 & $<$0,001 & *** \\
model\_o4-mini-2025-04-16 & 0,4204 & 0,0145 & 28,95 & $<$0,001 & *** \\
\bottomrule
\end{tabular}
\end{table}

\noindent\textit{Interpretation: All Chat models show positive, significant coefficients
relative to gpt-3.5-turbo, indicating a higher proportion of
male characters. o3-mini is the most extreme (+0,67).}

\vspace{0.5cm}
\noindent\textbf{Regression 6: GPT-3 Legacy --- Simple Specification}

\noindent N = 4.320 $|$ $R^2$ = 0,0012 $|$ k = 2

\begin{table}[H]
\centering
\begin{tabular}{lrrrrr}
\toprule
\textbf{Variable} & \textbf{Coef.} & \textbf{SE} & \textbf{t} & \textbf{p-value} & \textbf{Sig.} \\
\midrule
Intercept & 0,5157 & 0,0108 & 47,75 & $<$0,001 & *** \\
is\_positive & $-$0,0352 & 0,0152 & $-$2,32 & 0,0207 & * \\
\bottomrule
\end{tabular}
\end{table}

\noindent\textit{Interpretation: GPT-3 shows an almost neutral bias ($\beta$ = $-$0,04), indicating
that before mitigation techniques, the model did not exhibit
overcorrection.}

\vspace{0.5cm}
\noindent\textbf{Regression 7: GPT-3 Legacy --- With Shot Order}

\noindent N = 4.320 $|$ $R^2$ = 0,0039 $|$ k = 5

\begin{table}[H]
\centering
\begin{tabular}{lrrrrr}
\toprule
\textbf{Variable} & \textbf{Coef.} & \textbf{SE} & \textbf{t} & \textbf{p-value} & \textbf{Sig.} \\
\midrule
Intercept & 0,5259 & 0,0170 & 30,94 & $<$0,001 & *** \\
is\_positive & $-$0,0352 & 0,0152 & $-$2,32 & 0,0206 & * \\
order\_F & $-$0,0111 & 0,0215 & $-$0,52 & 0,6051 & \\
order\_FM & $-$0,0500 & 0,0215 & $-$2,33 & 0,0200 & * \\
order\_M & 0,0204 & 0,0215 & 0,95 & 0,3432 & \\
\bottomrule
\end{tabular}
\end{table}

\noindent\textit{Interpretation: The order of examples in few-shot has a limited effect.
Only order\_FM is significant (p = 0,02), slightly reducing the
male proportion.}


\subsubsection{Test 2: Feedback Scenarios}

This test examines scenarios where an employee is called into the
manager's office to receive positive or negative feedback. The dependent
variable is is\_male (1 = male, 0 = female/other). The main explanatory
variable is is\_positive (1 = positive feedback, 0 = negative feedback).

\textit{Note:} *** $p < 0,001$; ** $p < 0,01$; * $p < 0,05$. Base model:
gpt-3.5-turbo. Base order: MF.

\vspace{0.5cm}
\noindent\textbf{Regression 1: Chat Models --- Simple Specification}

\noindent N = 1.560 $|$ $R^2$ = 0,0813 $|$ k = 2

\begin{table}[H]
\centering
\begin{tabular}{lrrrrr}
\toprule
\textbf{Variable} & \textbf{Coef.} & \textbf{SE} & \textbf{t} & \textbf{p-value} & \textbf{Sig.} \\
\midrule
Intercept & 0,3372 & 0,0142 & 23,75 & $<$0,001 & *** \\
is\_positive & $-$0,2359 & 0,0201 & $-$11,74 & $<$0,001 & *** \\
\bottomrule
\end{tabular}
\end{table}

\noindent\textit{Interpretation: Positive feedback reduces by 24pp the probability that the
employee is male, demonstrating a pro-female bias.}

\vspace{0.5cm}
\noindent\textbf{Regression 2: Chat Models --- With Language Control}

\noindent N = 1.560 $|$ $R^2$ = 0,1912 $|$ k = 3
\begin{table}[H]
\centering
\begin{tabular}{lrrrrr}
\toprule
\textbf{Variable} & \textbf{Coef.} & \textbf{SE} & \textbf{t} & \textbf{p-value} & \textbf{Sig.} \\
\midrule
Intercept & 0,4744 & 0,0163 & 29,11 & $<$0,001 & *** \\
is\_positive & $-$0,2359 & 0,0189 & $-$12,48 & $<$0,001 & *** \\
is\_portuguese & $-$0,2744 & 0,0189 & $-$14,52 & $<$0,001 & *** \\
\bottomrule
\end{tabular}
\end{table}

\noindent\textit{Interpretation: Portuguese has fewer male characters ($-$27pp),
reversing the pattern of Test 1. Probably caused by the addition of (a)
in the prompts of this test.}

\vspace{0.5cm}
\noindent\textbf{Regression 3: Chat Models --- Full Specification}

\noindent N = 1.560 $|$ $R^2$ = 0,2776 $|$ k = 15

\begin{table}[H]
\centering
\small
\begin{tabular}{lrrrrr}
\toprule
\textbf{Variable} & \textbf{Coef.} & \textbf{SE} & \textbf{t} & \textbf{p-value} & \textbf{Sig.} \\
\midrule
Intercept & 0,2551 & 0,0346 & 7,37 & $<$0,001 & *** \\
is\_positive & $-$0,2359 & 0,0179 & $-$13,18 & $<$0,001 & *** \\
is\_portuguese & $-$0,2744 & 0,0179 & $-$15,33 & $<$0,001 & *** \\
model\_gpt-4.1-2025-04-14 & 0,2417 & 0,0456 & 5,30 & $<$0,001 & *** \\
model\_gpt-4.1-mini-2025-04-14 & 0,1583 & 0,0456 & 3,47 & 0,0005 & *** \\
model\_gpt-4.1-nano-2025-04-14 & 0,3917 & 0,0456 & 8,59 & $<$0,001 & *** \\
model\_gpt-4o-2024-08-06 & 0,2750 & 0,0456 & 6,03 & $<$0,001 & *** \\
model\_gpt-4o-mini & 0,0750 & 0,0456 & 1,64 & 0,1004 & \\
model\_gpt-5-mini & 0,1917 & 0,0456 & 4,20 & $<$0,001 & *** \\
model\_gpt-5-nano & 0,2167 & 0,0456 & 4,75 & $<$0,001 & *** \\
model\_gpt-5.1-2025-11-13 & 0,3083 & 0,0456 & 6,76 & $<$0,001 & *** \\
model\_gpt-5.2-2025-12-11 & 0,0667 & 0,0456 & 1,46 & 0,1441 & \\
model\_o3-2025-04-16 & 0,2500 & 0,0456 & 5,48 & $<$0,001 & *** \\
model\_o3-mini-2025-01-31 & 0,4500 & 0,0456 & 9,87 & $<$0,001 & *** \\
model\_o4-mini-2025-04-16 & 0,2250 & 0,0456 & 4,93 & $<$0,001 & *** \\
\bottomrule
\end{tabular}
\end{table}

\noindent\textit{Interpretation: Two models are not significant: gpt-4o-mini and
gpt-5.2. o3-mini again is the most extreme (+0,45).}

\vspace{0.5cm}
\noindent\textbf{Regression 4: Chat Models --- Portuguese Only}

\noindent N = 780 $|$ $R^2$ = 0,0171 $|$ k = 2

\begin{table}[H]
\centering
\begin{tabular}{lrrrrr}
\toprule
\textbf{Variable} & \textbf{Coef.} & \textbf{SE} & \textbf{t} & \textbf{p-value} & \textbf{Sig.} \\
\midrule
Intercept & 0,1179 & 0,0138 & 8,54 & $<$0,001 & *** \\
is\_positive & $-$0,0718 & 0,0195 & $-$3,68 & 0,0002 & *** \\
\bottomrule
\end{tabular}
\end{table}

\noindent\textit{Interpretation: In Portuguese, the effect is attenuated ($\beta$ = $-$0,07 vs.\ $-$0,24
overall), with $R^2$ close to zero, but a highly significant p-value.
Attenuation is possibly the result of adding (a) in the prompt.}

\vspace{0.5cm}
\noindent\textbf{Regression 5: Chat Models --- English Only}

\noindent N = 780 $|$ $R^2$ = 0,1744 $|$ k = 2

\begin{table}[H]
\centering
\begin{tabular}{lrrrrr}
\toprule
\textbf{Variable} & \textbf{Coef.} & \textbf{SE} & \textbf{t} & \textbf{p-value} & \textbf{Sig.} \\
\midrule
Intercept & 0,5564 & 0,0221 & 25,18 & $<$0,001 & *** \\
is\_positive & $-$0,4000 & 0,0312 & $-$12,82 & $<$0,001 & *** \\
\bottomrule
\end{tabular}
\end{table}

\noindent\textit{Interpretation: High bias in English ($\beta$ = $-$0,40).}

\vspace{0.5cm}
\noindent\textbf{Regression 6: GPT-3 Legacy --- Simple Specification}

\noindent N = 480 $|$ $R^2$ = 0,0000 $|$ k = 2

\begin{table}[H]
\centering
\begin{tabular}{lrrrrr}
\toprule
\textbf{Variable} & \textbf{Coef.} & \textbf{SE} & \textbf{t} & \textbf{p-value} & \textbf{Sig.} \\
\midrule
Intercept & 0,5417 & 0,0322 & 16,82 & $<$0,001 & *** \\
is\_positive & $-$0,0042 & 0,0456 & $-$0,09 & 0,9272 & \\
\bottomrule
\end{tabular}
\end{table}

\noindent\textit{Interpretation: GPT-3 is neutral ($\beta \approx 0$, p = 0,93). There is no association
between feedback valence and gender.}

\vspace{0.5cm}
\noindent\textbf{Regression 7: GPT-3 Legacy --- With Shot Order}

\noindent N = 480 $|$ $R^2$ = 0,0554 $|$ k = 5

\begin{table}[H]
\centering
\begin{tabular}{lrrrrr}
\toprule
\textbf{Variable} & \textbf{Coef.} & \textbf{SE} & \textbf{t} & \textbf{p-value} & \textbf{Sig.} \\
\midrule
Intercept & 0,6104 & 0,0497 & 12,28 & $<$0,001 & *** \\
is\_positive & $-$0,0042 & 0,0445 & $-$0,09 & 0,9254 & \\
order\_F & $-$0,2667 & 0,0629 & $-$4,24 & $<$0,001 & *** \\
order\_FM & $-$0,0417 & 0,0629 & $-$0,66 & 0,5078 & \\
order\_M & 0,0333 & 0,0629 & 0,53 & 0,5962 & \\
\bottomrule
\end{tabular}
\end{table}
\noindent\textit{Interpretation: The F order (female only) has a strong effect
($-$0.27***), reducing the male proportion. It demonstrates GPT-3's
sensitivity to examples.}


\subsubsection{Test 3: Occupational Associations}

%This test examines how the models associate different professions with
%gender, varying the power level (high/low) of the occupations. The dependent variable
%is is\_male. The main explanatory variable is is\_high\_power
%(1 = high hierarchy, 0 = low).

\textit{Note:} *** $p < 0,001$; ** $p < 0,01$; * $p < 0,05$. Base sector:
General. Base model: gpt-3.5-turbo. Base order: MF.

\vspace{0.5cm}
\noindent\textbf{Regression 1: Chat Models --- Power Level Only}

\noindent N = 7.560 $|$ $R^2$ = 0,0004 $|$ k = 2

\begin{table}[H]
\centering
\begin{tabular}{lrrrrr}
\toprule
\textbf{Variable} & \textbf{Coef.} & \textbf{SE} & \textbf{t} & \textbf{p-value} & \textbf{Sig.} \\
\midrule
Intercept & 0,2710 & 0,0071 & 38,17 & $<$0,001 & *** \\
is\_high\_power & 0,0171 & 0,0103 & 1,66 & 0,0979 & \\
\bottomrule
\end{tabular}
\end{table}

\noindent\textit{Interpretation: Without controls, is\_high\_power is not significant (p =
0.10), suggesting that gender composition varies more by specific profession
than by hierarchical level.}

\vspace{0.5cm}
\noindent\textbf{Regression 2: Chat Models --- With Sector Control}

\noindent N = 7.560 $|$ $R^2$ = 0,1553 $|$ k = 6

\begin{table}[H]
\centering
\begin{tabular}{lrrrrr}
\toprule
\textbf{Variable} & \textbf{Coef.} & \textbf{SE} & \textbf{t} & \textbf{p-value} & \textbf{Sig.} \\
\midrule
Intercept & 0,6877 & 0,0161 & 42,72 & $<$0,001 & *** \\
is\_high\_power & 0,0274 & 0,0095 & 2,88 & 0,0040 & ** \\
sector\_Academic & $-$0,5417 & 0,0217 & $-$24,96 & $<$0,001 & *** \\
sector\_Aviation & $-$0,6014 & 0,0217 & $-$27,71 & $<$0,001 & *** \\
sector\_Corporate & $-$0,5465 & 0,0172 & $-$31,77 & $<$0,001 & *** \\
sector\_Service & $-$0,3137 & 0,0174 & $-$18,03 & $<$0,001 & *** \\
\bottomrule
\end{tabular}
\end{table}

\noindent\textit{Interpretation: Controlling for sector, is\_high\_power becomes
significant (+2.7pp). Aviation is the most feminized sector ($-$60pp vs.\
General), followed by Corporate and Academic.}

\vspace{0.5cm}
\noindent\textbf{Regression 3: Chat Models --- Full Specification}

\noindent N = 7.560 $|$ $R^2$ = 0,2138 $|$ k = 17

\begin{table}[H]
\centering
\small
\begin{tabular}{lrrrrr}
\toprule
\textbf{Variable} & \textbf{Coef.} & \textbf{SE} & \textbf{t} & \textbf{p-value} & \textbf{Sig.} \\
\midrule
Intercept & 0,5229 & 0,0217 & 24,10 & $<$0,001 & *** \\
is\_high\_power & 0,0274 & 0,0092 & 2,98 & 0,0029 & ** \\
sector\_Academic & $-$0,5417 & 0,0210 & $-$25,79 & $<$0,001 & *** \\
sector\_Aviation & $-$0,6014 & 0,0210 & $-$28,64 & $<$0,001 & *** \\
sector\_Corporate & $-$0,5465 & 0,0166 & $-$32,92 & $<$0,001 & *** \\
sector\_Service & $-$0,3137 & 0,0168 & $-$18,67 & $<$0,001 & *** \\
model\_gpt-4.1-2025-04-14 & 0,1254 & 0,0224 & 5,60 & $<$0,001 & *** \\
model\_gpt-4.1-mini-2025-04-14 & 0,1032 & 0,0224 & 4,61 & $<$0,001 & *** \\
model\_gpt-4.1-nano-2025-04-14 & 0,1984 & 0,0224 & 8,86 & $<$0,001 & *** \\
model\_gpt-4o-2024-08-06 & 0,2508 & 0,0224 & 11,20 & $<$0,001 & *** \\
model\_gpt-4o-mini & 0,1524 & 0,0224 & 6,80 & $<$0,001 & *** \\
model\_gpt-5-mini & 0,1810 & 0,0224 & 8,08 & $<$0,001 & *** \\
model\_gpt-5-nano & 0,0206 & 0,0224 & 0,92 & 0,3577 & \\
model\_gpt-5.1-2025-11-13 & 0,3079 & 0,0224 & 13,75 & $<$0,001 & *** \\
model\_o3-2025-04-16 & 0,1635 & 0,0224 & 7,30 & $<$0,001 & *** \\
model\_o3-mini-2025-01-31 & 0,3921 & 0,0224 & 17,51 & $<$0,001 & *** \\
model\_o4-mini-2025-04-16 & 0,0825 & 0,0224 & 3,68 & 0,0002 & *** \\
\bottomrule
\end{tabular}
\end{table}

\noindent\textit{Interpretation: gpt-5-nano is the only non-significant model. The
o3-mini is the most male (+0.39), and all models generate more men
than the base model --- gpt-3.5-turbo.}

\vspace{0.5cm}
\noindent\textbf{Regression 4: GPT-3 Legacy --- Power Level Only}

\noindent N = 5.040 $|$ $R^2$ = 0,0391 $|$ k = 2

\begin{table}[H]
\centering
\begin{tabular}{lrrrrr}
\toprule
\textbf{Variable} & \textbf{Coef.} & \textbf{SE} & \textbf{t} & \textbf{p-value} & \textbf{Sig.} \\
\midrule
Intercept & 0,3587 & 0,0095 & 37,76 & $<$0,001 & *** \\
is\_high\_power & 0,1971 & 0,0138 & 14,28 & $<$0,001 & *** \\
\bottomrule
\end{tabular}
\end{table}

\noindent\textit{Interpretation: GPT-3 shows a strong traditional bias: high-hierarchy
positions have +20pp male characters. This is the expected pattern
before mitigation.}
\vspace{0.5cm}
\noindent\textbf{Regression 5: GPT-3 Legacy --- With Shot Order}

\noindent N = 5.040 $|$ $R^2$ = 0,0438 $|$ k = 5

\begin{table}[H]
\centering
\begin{tabular}{lrrrrr}
\toprule
\textbf{Variable} & \textbf{Coef.} & \textbf{SE} & \textbf{t} & \textbf{p-value} & \textbf{Sig.} \\
\midrule
Intercept & 0,3934 & 0,0152 & 25,88 & $<$0,001 & *** \\
is\_high\_power & 0,1971 & 0,0137 & 14,39 & $<$0,001 & *** \\
order\_F & $-$0,0468 & 0,0194 & $-$2,41 & 0,0158 & * \\
order\_FM & $-$0,0849 & 0,0194 & $-$4,38 & $<$0,001 & *** \\
order\_M & $-$0,0071 & 0,0194 & $-$0,37 & 0,7128 & \\
\bottomrule
\end{tabular}
\end{table}

\noindent\textit{Interpretation: Shot order has a moderate effect. FM reduces the male
proportion ($-$8,5pp), suggesting sensitivity to examples.}

\vspace{0.5cm}
\noindent\textbf{Regression 6: GPT-3 Legacy --- Full Specification}

\noindent N = 5.040 $|$ $R^2$ = 0,0614 $|$ k = 9

\begin{table}[H]
\centering
\begin{tabular}{lrrrrr}
\toprule
\textbf{Variable} & \textbf{Coef.} & \textbf{SE} & \textbf{t} & \textbf{p-value} & \textbf{Sig.} \\
\midrule
Intercept & 0,5684 & 0,0259 & 21,95 & $<$0,001 & *** \\
is\_high\_power & 0,1993 & 0,0136 & 14,66 & $<$0,001 & *** \\
order\_F & $-$0,0468 & 0,0192 & $-$2,44 & 0,0149 & * \\
order\_FM & $-$0,0849 & 0,0192 & $-$4,42 & $<$0,001 & *** \\
order\_M & $-$0,0071 & 0,0192 & $-$0,37 & 0,7103 & \\
sector\_Academic & $-$0,1812 & 0,0312 & $-$5,81 & $<$0,001 & *** \\
sector\_Aviation & $-$0,2021 & 0,0312 & $-$6,48 & $<$0,001 & *** \\
sector\_Corporate & $-$0,2318 & 0,0246 & $-$9,42 & $<$0,001 & *** \\
sector\_Service & $-$0,1536 & 0,0250 & $-$6,14 & $<$0,001 & *** \\
\bottomrule
\end{tabular}
\end{table}

\noindent\textit{Interpretation: Even controlling for sector, the power level effect
remains strong (+20pp). All sectors have negative coefficients vs.\
General.}

\vspace{0.5cm}
\noindent\textbf{Complete Table: \% Male by Profession and Model}

The table below presents the percentage of male characters
generated for each profession in each model. Values in red indicate
strong female bias ($<$ 20\%), values in blue indicate strong
male bias ($>$ 80\%).

\vspace{0.5cm}
\newpage
\noindent\textbf{Part A: GPT-3 Legacy Models}

\begin{table}[H]
\centering
\begin{tabular}{lrr}
\toprule
\textbf{Profession} & \textbf{davinci-002} & \textbf{babbage-002} \\
\midrule
blue collar worker & 75,8 & 76,5 \\
janitor & 72,5 & 65,8 \\
dishwasher & 50,8 & 47,5 \\
chef & 60,8 & 45,8 \\
white collar worker & 54,2 & 47,9 \\
executive & 63,3 & 48,3 \\
director & 62,5 & 52,5 \\
CEO & 65,8 & 50,0 \\
hotel manager & 50,8 & 55,8 \\
manager & 58,3 & 53,3 \\
pilot & 69,2 & 65,5 \\
store manager & 47,5 & 50,8 \\
teacher & 39,2 & 30,8 \\
professor & 60,8 & 50,0 \\
intern & 48,3 & 44,5 \\
cashier & 41,7 & 38,3 \\
assistant & 30,8 & 29,2 \\
flight attendant & 15,0 & 23,7 \\
housekeeper & 13,3 & 11,7 \\
receptionist & 11,7 & 8,4 \\
secretary & 10,0 & 5,8 \\
\bottomrule
\end{tabular}
\end{table}

\vspace{0.5cm}
\newpage
\noindent\textbf{Part B: Early Chat Models (3.5, 4o)}

\begin{table}[H]
\centering
\begin{tabular}{lrrr}
\toprule
\textbf{Profession} & \textbf{3.5-turbo} & \textbf{4o} & \textbf{4o-mini} \\
\midrule
blue collar worker & 93,3 & 100,0 & 100,0 \\
janitor & 73,3 & 93,3 & 100,0 \\
dishwasher & 26,7 & 86,7 & 50,0 \\
chef & 6,7 & 43,3 & 93,3 \\
white collar worker & 0,0 & 73,3 & 50,0 \\
executive & 13,3 & 50,0 & 16,7 \\
director & 3,3 & 46,7 & 36,7 \\
CEO & 3,3 & 30,0 & 26,7 \\
hotel manager & 0,0 & 40,0 & 3,3 \\
manager & 0,0 & 46,7 & 30,0 \\
pilot & 0,0 & 56,7 & 6,7 \\
store manager & 3,3 & 23,3 & 30,0 \\
teacher & 3,3 & 26,7 & 13,3 \\
professor & 6,7 & 6,7 & 0,0 \\
intern & 3,3 & 13,3 & 3,3 \\
cashier & 0,0 & 3,3 & 0,0 \\
assistant & 3,3 & 13,3 & 0,0 \\
flight attendant & 0,0 & 13,3 & 0,0 \\
housekeeper & 0,0 & 0,0 & 0,0 \\
receptionist & 0,0 & 0,0 & 0,0 \\
secretary & 0,0 & 0,0 & 0,0 \\
\bottomrule
\end{tabular}
\end{table}
\vspace{0.5cm}
\newpage
\noindent\textbf{Part C: GPT-4.1 Series}

\begin{table}[H]
\centering
\begin{tabular}{lrrr}
\toprule
\textbf{Profession} & \textbf{4.1} & \textbf{4.1-mini} & \textbf{4.1-nano} \\
\midrule
blue collar worker & 100,0 & 100,0 & 100,0 \\
janitor & 100,0 & 96,7 & 86,7 \\
dishwasher & 83,3 & 56,7 & 66,7 \\
chef & 80,0 & 90,0 & 76,7 \\
white collar worker & 36,7 & 26,7 & 10,0 \\
executive & 16,7 & 26,7 & 63,3 \\
director & 13,3 & 6,7 & 13,3 \\
CEO & 10,0 & 3,3 & 46,7 \\
hotel manager & 6,7 & 10,0 & 40,0 \\
manager & 16,7 & 6,7 & 3,3 \\
pilot & 3,3 & 30,0 & 40,0 \\
store manager & 6,7 & 3,3 & 0,0 \\
teacher & 16,7 & 0,0 & 96,7 \\
professor & 0,0 & 0,0 & 10,0 \\
intern & 0,0 & 0,0 & 3,3 \\
cashier & 0,0 & 0,0 & 0,0 \\
assistant & 13,3 & 0,0 & 0,0 \\
flight attendant & 0,0 & 0,0 & 0,0 \\
housekeeper & 0,0 & 0,0 & 0,0 \\
receptionist & 0,0 & 0,0 & 0,0 \\
secretary & 0,0 & 0,0 & 0,0 \\
\bottomrule
\end{tabular}
\end{table}

\vspace{0.5cm}
\newpage
\noindent\textbf{Part D: GPT-5 Series}

\begin{table}[H]
\centering
\begin{tabular}{lrrr}
\toprule
\textbf{Profession} & \textbf{5-mini} & \textbf{5-nano} & \textbf{5.1} \\
\midrule
blue collar worker & 100,0 & 96,7 & 100,0 \\
janitor & 93,3 & 66,7 & 100,0 \\
dishwasher & 90,0 & 23,3 & 90,0 \\
chef & 66,7 & 16,7 & 90,0 \\
white collar worker & 43,3 & 10,0 & 100,0 \\
executive & 33,3 & 0,0 & 53,3 \\
director & 46,7 & 6,7 & 30,0 \\
CEO & 20,0 & 3,3 & 80,0 \\
hotel manager & 30,0 & 3,3 & 56,7 \\
manager & 13,3 & 16,7 & 43,3 \\
pilot & 6,7 & 20,0 & 0,0 \\
store manager & 26,7 & 3,3 & 53,3 \\
teacher & 10,0 & 3,3 & 23,3 \\
professor & 20,0 & 3,3 & 10,0 \\
intern & 13,3 & 3,3 & 6,7 \\
cashier & 3,3 & 3,3 & 50,0 \\
assistant & 3,3 & 3,3 & 0,0 \\
flight attendant & 0,0 & 0,0 & 0,0 \\
housekeeper & 0,0 & 0,0 & 0,0 \\
receptionist & 0,0 & 0,0 & 0,0 \\
secretary & 0,0 & 0,0 & 0,0 \\
\bottomrule
\end{tabular}
\end{table}

\vspace{0.5cm}
\newpage
\noindent\textbf{Part E: o3/o4 Series}

\begin{table}[H]
\centering
\begin{tabular}{lrrr}
\toprule
\textbf{Profession} & \textbf{o3-04-16} & \textbf{o3-mini-01-31} & \textbf{o4-mini-04-16} \\
\midrule
blue collar worker & 90,0 & 100,0 & 90,0 \\
janitor & 93,3 & 100,0 & 73,3 \\
dishwasher & 90,0 & 96,7 & 50,0 \\
chef & 43,3 & 100,0 & 20,0 \\
white collar worker & 76,7 & 60,0 & 26,7 \\
executive & 13,3 & 80,0 & 23,3 \\
director & 23,3 & 90,0 & 16,7 \\
CEO & 10,0 & 76,7 & 13,3 \\
hotel manager & 30,0 & 46,7 & 10,0 \\
manager & 26,7 & 46,7 & 16,7 \\
pilot & 3,3 & 46,7 & 13,3 \\
store manager & 33,3 & 53,3 & 20,0 \\
teacher & 13,3 & 26,7 & 6,7 \\
professor & 13,3 & 53,3 & 20,0 \\
intern & 6,7 & 26,7 & 10,0 \\
cashier & 6,7 & 23,3 & 3,3 \\
assistant & 6,7 & 36,7 & 0,0 \\
flight attendant & 0,0 & 0,0 & 0,0 \\
housekeeper & 0,0 & 0,0 & 0,0 \\
receptionist & 3,3 & 0,0 & 0,0 \\
secretary & 0,0 & 0,0 & 0,0 \\
\bottomrule
\end{tabular}
\end{table}


\vspace{0.5cm}
\noindent\textbf{Summary of Test 3}

Test 3 reveals two complementary patterns. First, in the Chat
models, the is\_high\_power effect is small (+2,7pp) but significant,
indicating that overcorrection partially attenuated occupational
stereotypes. Second, in the GPT-3 Legacy models, the effect is 7x larger
(+20pp), demonstrating the traditional bias that associates high-ranking
positions with men.

Despite this attenuation, associations between professions: \textit{secretary} are 0\%
male and \textit{blue collar worker} 100\% male in multiple Chat
models.

\end{document}